\section{Implementation: Classic Visitor in C++}
\label{sec:implementation}

To demonstrate the practical application of the theoretical concepts discussed, this section presents a complete, runnable C++ implementation of the Classic Visitor pattern applied to the geometric shapes scenario \cite{gof1994}.

In a professional C++ software architecture, these components are typically distributed across multiple header and source files to minimize compilation dependencies. A critical technical nuance when implementing the Visitor pattern in C++ is the inherently circular dependency between the \texttt{Visitor} interface and the \texttt{Element} interface. To resolve this, forward declarations must be utilized appropriately. The implementation is logically divided into three primary components: core interfaces, concrete elements and operations, and client execution.

\subsection{Core Interfaces}
The foundational step involves defining the abstract base classes. Forward declarations of the concrete elements are required before defining the \texttt{Visitor} interface, allowing the compiler to recognize the types without requiring their full definitions.

\begin{lstlisting}[language=C++, caption=Core Interfaces (typically \texttt{Visitor.h} and \texttt{Shape.h})]
#pragma once
#include <string>

// Forward declarations to resolve circular dependencies
class Circle;
class Rectangle;

// ================= VISITOR HIERARCHY =================
class Visitor {
public:
    virtual ~Visitor() = default;
    
    // Overloaded visit methods for each Concrete Element
    virtual void visit(Circle& c) = 0;
    virtual void visit(Rectangle& r) = 0;
};

// ================= ELEMENT HIERARCHY =================
class Shape {
public:
    virtual ~Shape() = default;
    virtual void draw() = 0;
    virtual double area() = 0;
    
    // The core entry point for the Visitor
    virtual void accept(Visitor& v) = 0;
};
\end{lstlisting}

\subsection{Concrete Implementations}
Once the abstract contracts are established, the concrete classes can be implemented. The concrete elements (\texttt{Circle}, \texttt{Rectangle}) provide their domain-specific logic and implement the \texttt{accept} method to initiate the second dispatch. Concurrently, the \texttt{SaveVisitor} encapsulates the state (e.g., the file path) and the specific serialization logic.

\begin{lstlisting}[language=C++, caption=Concrete Classes (Elements and Visitors)]
#include <iostream>
// In a real project, this file would include Shape.h and Visitor.h

// --- Concrete Elements ---
class Circle : public Shape {
private:
    float radius;
public:
    explicit Circle(float r) : radius(r) {}
    
    void draw() override { std::cout << "Drawing Circle\n"; }
    double area() override { return 3.14159 * radius * radius; }
    
    // First Dispatch: Resolves to Circle::accept
    void accept(Visitor& v) override {
        // Second Dispatch: Calls visit(Circle&)
        v.visit(*this); 
    }
};

class Rectangle : public Shape {
private:
    float a, b;
public:
    Rectangle(float width, float height) : a(width), b(height) {}
    
    void draw() override { std::cout << "Drawing Rectangle\n"; }
    double area() override { return a * b; }
    
    // First Dispatch: Resolves to Rectangle::accept
    void accept(Visitor& v) override {
        // Second Dispatch: Calls visit(Rectangle&)
        v.visit(*this); 
    }
};

// --- Concrete Visitor ---
class SaveVisitor : public Visitor {
private:
    std::string filePath;
public:
    // Explicit constructor to initialize the visitor's state
    explicit SaveVisitor(const std::string& path) : filePath(path) {}
    
    void visit(Circle& c) override {
        // Logic for saving a Circle to file
        std::cout << "Saving Circle data to " << filePath << "...\n";
    }
    
    void visit(Rectangle& r) override {
        // Logic for saving a Rectangle to file
        std::cout << "Saving Rectangle data to " << filePath << "...\n";
    }
};
\end{lstlisting}

\subsection{Client Code Execution}
The client code remains entirely decoupled from the concrete visitor logic. It simply iterates through a heterogeneous collection of shapes and applies the desired operation via the \texttt{accept} method, showcasing the elegance and flexibility of the architecture.

\begin{lstlisting}[language=C++, caption=Client Code demonstrating polymorphic iteration and Double Dispatch (\texttt{main.cpp})]
#include <vector>
#include <memory>

void practiceUsage() {
    // Constructing a heterogeneous collection of shapes using smart pointers for RAII
    std::vector<std::unique_ptr<Shape>> shapes;
    shapes.push_back(std::make_unique<Circle>(5.0f));
    shapes.push_back(std::make_unique<Rectangle>(4.0f, 6.0f));
    
    // Instantiating the new operation with necessary state
    SaveVisitor saveOp("export_data.xml");
    
    // Applying the operation across the entire object structure
    for (const auto& shape : shapes) {
        // Executes double dispatch
        shape->accept(saveOp); 
    }
    
    // No explicit memory cleanup required (handled by std::unique_ptr)
}
\end{lstlisting}
